Tabelle sollten grundsätzlich so wenig Trennlinien wie möglich enthalten. Auf
vertikale Linien sollte man komplett verzichten. Horizontale Linien sollte man
nur zur Trennung des Tabellenkopfs oder -fußes einziehen. Die
\LaTeX-Standardkommandos sind dahingehend nicht optimal. Wir raten zur
Verwendung des Pakets~\ctanref{booktabs}.

\begin{table}
\centering
\begin{tabular}{lcr}
\toprule
Spalte 1 & Spalte 2 & Spalte 3 \\
\midrule
links    & mitte    & rechts   \\
l        & m        & r        \\
\bottomrule
\end{tabular}

\caption{Tabelle mit \ctanpkg{booktabs}}
\label{tab:latex:tab:booktabs}
\end{table}

\lstinputlisting[language={[LaTeX]{TeX}},style=block,float,caption={Tabelle mit \ctanpkg{booktabs}},label={tab:latex:tab:booktabs}]{code/latex_tab_booktabs.tex}

Die Trennlinien in Tabelle~\ref{tab:latex:tab:booktabs} wurden mit dem
Paket~\ctanpkg{booktabs} gesetzt. Der Code dafür ist in
Listing~\ref{tab:latex:tab:booktabs} abgedruckt. Neben den Befehlen
\lstinline[language={[LaTeX]{TeX}}]+\toprule+,
\lstinline[language={[LaTeX]{TeX}}]+\midrule+ und
\lstinline[language={[LaTeX]{TeX}}]+\bottomrule+ gibt es noch weitere Kommandos
zu Trennlinien, die sich nur über bestimmte Spalten erstrecken, oder den
Zeilenabstands anpassen. Hierfür sei auf die Paketdokumentation verwiesen.
